\begin{abstract}
1848年，Bezzel提出八皇后问题：在$8\times8$棋盘上摆八个皇后，要求彼此互不攻击——
即任意两个不同行、不同列、不同对角线。这后来被推广为N皇后问题：
在$N\times N$棋盘上放$N$个互不攻击的皇后，问有多少种放法，记作$Q(N)$。
$Q(N)$随$N$超指数增长，穷举搜索到2025年也只算到$N=27$。
2022年，Simkin严格证明了渐近式$Q(N)\sim(N/e^\gamma)^N$，随后Nobel等人
用大规模Newton方法把常数压到了$\gamma=1.94400(1)$。
不过，此前所有对$\gamma$的测定走的都是组合数学的路子，缺少物理学的交叉验证。

我们把N皇后问题转化成一个格点气体：每个皇后是一颗占据粒子，凡在同一行、
同一列、同一对角线上的两颗粒子，就产生一个单位的排斥能。这样一来，
哈密顿量$H=J\sum n_{ij}n_{i'j'}$恰好统计了所有互攻的皇后对数，
基态$E=0$正好就是$Q(N)$种非攻击解。从每皇后的基态熵
$s_0=\ln N-\gamma$可以读出清晰的约束层级：$N\to\infty$时，
禁止列重复吃掉正好$1$单位熵（来自Stirling公式$\ln N!=N\ln N-N$），
两族对角线再各吃掉约$0.47$单位，合起来就是$\gamma-1\approx0.94$。
我们还导出了高温能量对任意有限$N$的严格公式
$E/N|_{T\to\infty}=(5N-1)(N-1)/[3N(N+1)]$，大$N$极限为$5/3$。

模拟方面，我们在八个尺寸（$N=8$到$1024$，逐次翻倍）上跑了正则系综Monte Carlo，
温度覆盖$T=0.05\sim500\,J$、等距布了$280$个点，每个点扫$10^8$步，
用的是Kawasaki粒子-空位交换加Metropolis判据。
结果显示，$N\ge32$以后每皇后比热$C_v/N$就坍塌到了同一条与$N$无关的曲线上，
峰值$C_v^{\max}/N\approx1.63$出现在$T^*\approx0.235\,J$。
从$N=32$到$1024$，峰高只变了$0.007$——如果有对数发散，按$\ln N$估算该有
$\sim0.06$/倍，实测远小于此，所以这不是相变，是个Schottky型异常。

有了$C_v/N$的收敛性和高温熵的精确值$S(\infty)/N=\ln N+1$，就可以对
$\int_0^\infty C_v/(NT)\,\mathrm{d}T$做热力学积分，从纯模拟数据里把Simkin常数提出来。
在$N=1024$处得到$\gamma_{\rm MC}=1.946\pm0.003$，跟组合学结果$1.94400(1)$只差
$0.11\%$（约$0.002$），落在jackknife误差$\pm0.003$以内。
这等于用物理手段独立检验了$\gamma$的值。
$\gamma_{\rm MC}(N)$随$N$单调靠近$\gamma$的趋势表明，剩下的那点偏差应该来自
有限尺寸修正。做个简单的标度外推$\gamma_{\rm MC}(N)=\gamma_\infty+aN^{-\alpha}$，
得到$\gamma_\infty=1.947\pm0.004$。

作为与Monte Carlo互补的另一条路线，我们还搭了一套基于矩阵乘积算符的张量网络
精确计数方案。四个方向的攻击约束被压进bond dimension $D=2$的MPO里，
交汇后产生只含$17$个非零元的局域张量，用边界MPS逐行收缩就能直接算出$Q(N)$。
有意思的是，该转移矩阵是幂零的（$\mathcal{T}^{N+1}=0$），这意味着DMRG、
CTMRG、VUMPS这些依赖最大特征值的标准谱方法在这里全都不适用。
本文为Simkin常数提供了一条完全独立于组合枚举的热力学测量途径。
\end{abstract}
